\ssr{ЗАКЛЮЧЕНИЕ}


В ходе выполнения лабораторной работы поставленная цель была достигнута: 
разработано программное обеспечение для валидации данных PDF-документов и проведён сравнительный анализ эффективности трёх больших языковых моделей при решении инженерных задач.

Таким образом, поставленные задачи были выполнены:
\begin{enumerate}
    \item выполнен аналитический обзор современных БЯМ, обоснован выбор моделей Gemini 3 Pro (SOTA), YandexGPT 5.1 (отечественная) и Qwen 3 Coder (локальная);
    \item спроектированы и отлажены регулярные выражения для выявления ошибок вёрстки (отсутствия неразрывного пробела) в специальных сочетаниях слов и чисел с учётом специфики Unicode;
    \item реализован алгоритм поиска подстрок в PDF-файле. Для решения проблемы отсутствия нативной нумерации строк в библиотеке \texttt{PyPDF2} применён метод динамического вычисления координат через анализ символов переноса в извлечённом тексте;
    \item разработано и протестировано консольное программное обеспечение, принимающее на вход путь к файлу и возвращающее детализированный отчёт о найденных несоответствиях;
    \item задокументирован процесс взаимодействия с моделями (Prompt Engineering), зафиксированы все итерации: от первичных ошибок (галлюцинации с библиотеками, ошибки области видимости) до получения рабочего кода;
    \item проведена проверка реализации на тестовом наборе файлов, результаты сведены в таблицу, подтверждающую корректность работы итогового алгоритма;
    \item выполнен исследовательский анализ стабильности генерации, который показал неочевидные результаты.
    \begin{itemize}
        \item режим рассуждений (Reasoning) у модели Gemini 3 Pro обеспечил наилучшее понимание бизнес-логики;
        \item для модели YandexGPT 5.1 отключение режима рассуждений, напротив, повысило качество кода, устранив логические ошибки;
        \item локальная модель Qwen 3 Coder продемонстрировала высокую техническую грамотность, но склонность к инверсии условий задачи.
    \end{itemize}
\end{enumerate}

\clearpage